\documentclass[11pt]{article}

\usepackage{amsmath,amssymb,amsthm}
\usepackage[margin=1in]{geometry}
\usepackage{hyperref}
\usepackage{enumitem}

\title{Audit of the CoHA-to-$\mathcal{W}_{1+\infty}$ Treatise:
Inconsistencies, Corrections, and the Correct Dictionary}
\author{Raeez Lorgat}
\date{2026-04-22}

\newtheorem{theorem}{Theorem}[section]
\newtheorem{proposition}[theorem]{Proposition}
\newtheorem{lemma}[theorem]{Lemma}
\newtheorem{corollary}[theorem]{Corollary}
\theoremstyle{definition}
\newtheorem{definition}[theorem]{Definition}
\newtheorem{example}[theorem]{Example}
\theoremstyle{remark}
\newtheorem{remark}[theorem]{Remark}
\newtheorem{flaw}[theorem]{Flaw}
\newtheorem{correction}[theorem]{Correction}

\newcommand{\CoHA}{\mathrm{CoHA}}
\newcommand{\Winf}{\mathcal{W}_{1+\infty}}
\newcommand{\hCS}{\mathrm{hCS}}
\newcommand{\HolFA}{\mathrm{HolFA}}
\newcommand{\Obs}{\mathrm{Obs}}
\newcommand{\CE}{\mathrm{CE}}
\newcommand{\Conf}{\mathrm{Conf}}
\newcommand{\Hilb}{\mathrm{Hilb}}
\newcommand{\Rep}{\mathrm{Rep}}
\newcommand{\Sym}{\mathrm{Sym}}
\newcommand{\gghone}{\widehat{\mathfrak{gl}}_1}
\newcommand{\ggtwo}{\widehat{\widehat{\mathfrak{gl}}_2}}
\newcommand{\ZZ}{\mathbb{Z}}
\newcommand{\CC}{\mathbb{C}}
\newcommand{\RR}{\mathbb{R}}
\newcommand{\NN}{\mathbb{N}}
\newcommand{\QQ}{\mathbb{Q}}
\newcommand{\kch}{\kappa_{\mathrm{ch}}}
\newcommand{\kcat}{\kappa_{\mathrm{cat}}}
\newcommand{\kBKM}{\kappa_{\mathrm{BKM}}}
\newcommand{\kfib}{\kappa_{\mathrm{fiber}}}
\newcommand{\Lie}{\mathrm{Lie}}
\newcommand{\HH}{\mathrm{HH}}
\newcommand{\ch}{\mathrm{ch}}
\newcommand{\ev}{\mathrm{ev}}
\newcommand{\cZ}{\mathcal{Z}}
\newcommand{\cF}{\mathcal{F}}
\newcommand{\cE}{\mathcal{E}}
\newcommand{\cO}{\mathcal{O}}
\newcommand{\cM}{\mathcal{M}}
\newcommand{\cD}{\mathcal{D}}
\newcommand{\cW}{\mathcal{W}}
\newcommand{\fg}{\mathfrak{g}}
\newcommand{\fgl}{\mathfrak{gl}}
\newcommand{\fsl}{\mathfrak{sl}}
\newcommand{\fh}{\mathfrak{h}}
\newcommand{\ftheis}{\mathfrak{heis}}
\newcommand{\Stone}{E_1}
\newcommand{\Stwo}{E_2}
\newcommand{\Sthree}{E_3}
\newcommand{\Eone}{E_1}
\newcommand{\Etwo}{E_2}
\newcommand{\Ethree}{E_3}

\hypersetup{colorlinks=true,linkcolor=black,citecolor=black,urlcolor=black}

\begin{document}

\maketitle

\begin{abstract}
The treatise \texttt{notes/CoHA\_to\_W\_infty\_treatise.tex} develops
three worked examples of the CoHA-to-vertex-algebra-to-gauge-theory
chain (affine three-space, the resolved conifold, the product
$K3 \times E$). This note audits the treatise sentence by sentence
against the Vol III programme crystallised in
\texttt{working\_notes.tex},
\texttt{notes/platonic\_synthesis\_waves\_11\_through\_16.tex}, and
\texttt{appendices/first\_principles\_cache.md}. Five attack-heal
cycles produce a catalogue of inconsistencies --- cache-violating
identifications, mis-stated theorem attributions, scope conflations,
symbol-definition gaps --- each paired with a rigorous correction and
anchored at the primary-literature strength the cache already records.
The consolidated correct dictionary positions $\Winf$ as the image of
$Y^+(\gghone)$ under evaluation at a one-parameter slice of the
truncation parameter, not as an identification, with the two algebras
distinguished by associativity class (CoHA: $\Eone$) and by the
chirality generated by the $\Omega$-background on the transverse
complex surface. The audit log records nine structural corrections,
six terminological or scope refinements, and four open questions
sharpened by the audit.
\end{abstract}

\section*{Conventions}

$\kappa_\bullet$ always carries a subscript: $\kch$ for the
CY-categorical chiral invariant computed via $\Phi$, $\kcat$ for the
Euler characteristic $\chi(\cO_X)$, $\kBKM$ for the Borcherds weight,
$\kfib$ for the fibre-lattice contribution. Bare $\kappa$ is forbidden
(Pattern AP113). $\Phi$ is the CY-to-chiral functor: ONE output per
CY category; $d$-dependent ($\Etwo$-chiral at $d \leq 2$,
$\Eone$-chiral at $d \geq 3$). Epsilon weights: $(\epsilon_1,
\epsilon_2, \epsilon_3)$ on $\CC^3$ with the $\Omega$-background
parameters on the transverse $\CC^2$ being $(\epsilon_1, \epsilon_2)$
and the Calabi--Yau identity $\epsilon_1 + \epsilon_2 + \epsilon_3 = 0$
written with $\epsilon_3 = -\epsilon_1 - \epsilon_2$. Primary citations
are retained in the cache-canonical form (Schiffmann--Vasserot 2013;
Tsymbaliuk 2017; Gaiotto--Rap\v{c}\'ak 2017; Costello 2013;
Davison--Meinhardt 2020; Gritsenko--Nikulin 1998; Oberdieck--Pixton
2017; Costello--Gwilliam 2017; Prochazka--Rap\v{c}\'ak 2018).

\tableofcontents

\section{Reading of the treatise against the programme}
\label{sec:reading}

The treatise advances three worked examples --- $\CC^3$, resolved
conifold, $K3 \times E$ --- as successive difficulty steps in the
chain
\[
\text{CY}_3 \;\to\; \text{CoHA} \;\to\; \text{vertex algebra}
\;\to\; \text{gauge theory}.
\]
Each example is staged through the same modules: quiver with
potential, Jacobi algebra, moduli stack, CoHA multiplication, positive
half, Drinfeld double, vertex-algebra image, gauge-theory origin,
holomorphic Chern--Simons descent. The treatise treats the
modules as parallel: a clean presentation at $\CC^3$, a partial
presentation at the conifold, and an enumeration of obstructions at
$K3 \times E$.

Directionally the treatise is correct on each of the
module-by-module progressions, and several blocks (the Jordan triple
Jacobi computation; the Tsymbaliuk Drinfeld double; the Szendroi
non-commutative crepant resolution of the conifold; the
Oberdieck--Pixton match to $\Phi_{10}$ on the primitive K3 class)
match the programme's recorded primary-literature strengths exactly.
Its structural problems are concentrated in four places: (a) the
framing of $\CoHA(\CC^3) \leftrightarrow \Winf$, (b) the handling of
the $\Omega$-background parameters and the CY$_3$ identity, (c) the
scope and attribution of the conifold and $K3 \times E$ CoHA
identifications, and (d) six-dimensional hCS-to-chiral output
containment. These are the foci of Section~\ref{sec:inconsistency}.

\section{Inconsistency catalogue with corrections}
\label{sec:inconsistency}

The audit is organised by descending structural load. Each item is
given as flaw (treatise text or implication), correction (rigorous
statement with primary-literature anchor), and a short remark on
why the flaw was structurally attractive.

\subsection{$\CoHA(\CC^3)$ is \emph{not} $\Winf$}
\label{subsec:coha-vs-winf}

\begin{flaw}[Identification of the CoHA with $\Winf$ as an
associative algebra]
\label{flaw:coha-winf}
The treatise moves, at Subsection~``Drinfeld double and identification
with $\Winf$-affine Yangian'', from
$\CoHA(\CC^3) \cong Y^+(\gghone)$ (Schiffmann--Vasserot 2013) to
$Y(\gghone)$ as a Drinfeld double and then to $\Winf[\lambda]$ via
``state-field correspondence on the vacuum module''. The summary
tabular row ``Identified with VOA: $\Winf$'' in the $\CC^3$ column
reads the chain as an identification between $\CoHA(\CC^3)$ and
$\Winf$. The cache-canonical statement is that they are two
different algebraic structures: $\CoHA(\CC^3) = Y^+(\gghone)$ is
associative ($\Eone$, no inner OPE), while $\Winf$ is a vertex
operator algebra on a curve.
\end{flaw}

\begin{correction}[The precise relation between
$Y^+(\gghone)$ and $\Winf$]
\label{cor:coha-winf}
The three algebras $Y^+(\gghone)$, $Y(\gghone)$, and $\Winf[\lambda]$
are related by the following strictly typed maps, none of which is an
identification:
\begin{enumerate}[label=\textup{(\roman*)}]
\item Schiffmann--Vasserot 2013, Thm.~1.1:
$\CoHA(\CC^3) \otimes_\FF \FF(\!(z)\!) \cong \mathrm{Sh}_{\gghone}$,
an associative shuffle algebra with kernel
$\omega(z,w) = (z-w-\epsilon_1)(z-w-\epsilon_2)(z-w-\epsilon_3)/(z-w)^3$.
\item Tsymbaliuk 2017, Thm.~1.1: the Drinfeld double
$Y(\gghone) := Y^+ \bowtie Y^0 \bowtie Y^-$ is topologically
Hopf-isomorphic to the affine Yangian of $\gghone$ in Drinfeld-currents
presentation.
\item Prochazka--Rap\v{c}\'ak 2018; Gaiotto--Rap\v{c}\'ak 2017
(\texttt{arXiv:1703.00982}): there is a one-parameter family
$Y(\gghone) \to \mathrm{End}(\text{vacuum module of } \Winf[\lambda])$
via the state-field correspondence on Fock modules, making $\Winf$ a
representation-theoretic target of the affine Yangian rather than an
associative-algebra identification.
\end{enumerate}
In the Vol III language:
$\CoHA(\CC^3) = Y^+(\gghone)$ is the positive half of the affine
Yangian, an $\Eone$ object; the full affine Yangian $Y(\gghone)$ acts
on $\Winf$; $\Winf$ itself is a $\Etwo$-chiral algebra in its own right
(CFT on a curve), and is recovered from $Y(\gghone)$ only after
evaluation at a slice of the truncation parameter together with a
completion. Symbolically,
\[
\CoHA(\CC^3) \;=\; Y^+(\gghone) \;\neq\; \Winf,
\quad
Y^+(\gghone) \hookrightarrow Y(\gghone)
\xrightarrow{\;\ev_\lambda\;} \mathrm{End}(\Winf[\lambda]),
\]
with $\ev_\lambda$ a representation, not an isomorphism of algebras
(and certainly not an isomorphism of associativity class).
\end{correction}

\begin{remark}
The $\CC^3$ Rosetta section of \texttt{working\_notes.tex}
(lines 1020--1166) displays this precisely: the functor-chain arrow
from $\Rep^{\Eone}(Y^+(\gghone))$ to
$\Rep^{\Etwo}(Y(\gghone)) \simeq \Rep^{\Etwo}(\Winf)$ goes through the
Drinfeld center $\cZ$, not through a direct identification. The
``five independent routes to $\Winf$'' subsection (lines 1658--1674)
is at the level of $\Winf$ itself as an $\Etwo$-chiral algebra, with
$\CoHA(\CC^3) = Y^+(\gghone)$ appearing as route~(a), not as $\Winf$
itself.
\end{remark}

\begin{correction}[$\Winf$ is the evaluation slice, not a completion
or deformation of $Y^+$]
\label{cor:eval-slice}
$\Winf$ is neither a completion of $Y^+(\gghone)$ nor an $\hbar$-scale
deformation. It is a representation-theoretic image: under the
Gaiotto--Rap\v{c}\'ak $\lambda$-parameter, one obtains the
endomorphism vertex algebra of the vacuum module of $Y(\gghone)$,
and this endomorphism vertex algebra is $\Winf[\lambda]$. The
$Y^+$-to-$\Winf$ map is a composition
\[
Y^+(\gghone) \hookrightarrow Y(\gghone) \xrightarrow{\;\ev_\lambda\;}
\mathrm{End}(\Winf[\lambda])\text{-vacuum},
\]
and the treatise's $\CC^3$ summary row would be accurate if it read
``$Y^+ \hookrightarrow Y \to \Winf$-endomorphisms'' rather than
``identified with VOA''.
\end{correction}

\subsection{The $\lambda_{\mathrm{Tr}}$-versus-$\lambda_{\mathrm{GR}}$
parametrisation and CY$_3$}
\label{subsec:lambda-gr-tr}

\begin{flaw}[$\lambda$ at CY$_3$ point is \emph{not} a number]
\label{flaw:lambda-cy3}
The ``$\lambda_{\mathrm{Tr}}/\lambda_{\mathrm{GR}}$'' remark and the
accompanying correction paragraph are largely correct, but partially
miswritten: the sentence
\emph{``only at the further special point $\epsilon_1 = \epsilon_2$
(i.e.\ $\epsilon_1 + \epsilon_2 = -\epsilon_3$ with $\epsilon_1
= \epsilon_2$) does $\lambda_{\mathrm{Tr}}$ collapse to a particular
numerical value''} asserts a diagonal slice as the only special point.
Two special points matter: the self-dual slice $\epsilon_1
+ \epsilon_2 = 0$ (gives $\Winf$ at $c = 1$, the Heisenberg VOA $H_1$;
see \texttt{working\_notes.tex} Thm.~\texttt{wn:thm:kappa-c3}), and the
``unrefined'' slice $\epsilon_1 = \epsilon_2$. The treatise conflates
them.
\end{flaw}

\begin{correction}[Two distinguished slices on the CY$_3$ parameter
surface]
\label{cor:lambda-slices}
After imposing $\epsilon_1 + \epsilon_2 + \epsilon_3 = 0$, the
$\Winf[\lambda]$ family depends on a single ratio. Two
representation-theoretically distinguished points (not a single one)
are:
\begin{itemize}
\item \emph{Self-dual / Heisenberg point} $\epsilon_1 + \epsilon_2 =
0$, equivalently $\epsilon_3 = 0$. Here $\Winf[\lambda]|_{\lambda \to
\infty}$ truncates to the Heisenberg VOA $H_1$ at $c = 1$ (see
\texttt{working\_notes.tex} line~1045: ``At the self-dual level
($\sigma_3 \to 0$, giving $h_1 = 1, h_2 = 0, h_3 = -1$), the output at
Step~3 is $\Winf$ at $c = 1$, which is the Heisenberg VOA $H_1$.'').
\item \emph{Unrefined point} $\epsilon_1 = \epsilon_2$: the Miki
$S_3$-symmetric locus after the CY$_3$ constraint, with
$\epsilon_3 = -2\epsilon_1$.
\end{itemize}
Neither slice is ``$\lambda = -1$''; the historical ``slip'' the
treatise alludes to is precisely the conflation of the self-dual
slice with the Gaiotto--Rap\v{c}\'ak evaluation at a numerical
$\lambda$.
\end{correction}

\subsection{Calabi--Yau condition: CY$_3$ identity on parameters vs.
CY$_3$ condition on $J(Q, W)$}
\label{subsec:cy3-identity}

\begin{flaw}[Treatise equates CY$_3$ parameter identity with CY$_3$
algebra condition]
\label{flaw:cy3-identity-conflation}
The paragraph starting ``$T = (\CC^\times)^3$ acts on $(X, Y, Z)$ with
weights $(t_1, t_2, t_3)$. The CY$_3$ condition (coming from the
holomorphic volume form $dX \wedge dY \wedge dZ$ being $T$-invariant)
is $t_1 t_2 t_3 = 1$'' packs two distinct Calabi--Yau conditions into
one statement. The CY$_3$ condition on the algebra $J(Q, W) =
k[X, Y, Z]$ is that the bimodule resolution is symmetric of length~3
(Ginzburg 2006; Bocklandt 2008); the $T$-weight identity $t_1 t_2
t_3 = 1$ is one specific realisation via the holomorphic volume form
on the transverse $\CC^3$, which is not the full CY$_3$ algebraic
datum.
\end{flaw}

\begin{correction}[Three compatible CY$_3$ conditions on $\CC^3$]
\label{cor:three-cy3}
Three a-priori different conditions agree on $\CC^3$ and should be
kept distinct in the presentation:
\begin{enumerate}[label=\textup{(\roman*)}]
\item \emph{Algebraic CY$_3$ condition on $J(Q, W)$}: existence of a
4-periodic, symmetric, length-3 self-dual bimodule resolution
(Ginzburg 2006 Def.~3.6.1; Bocklandt 2008).
\item \emph{Geometric CY$_3$ condition on $X = \CC^3$}: holomorphic
volume form $\Omega = dz_1 \wedge dz_2 \wedge dz_3$.
\item \emph{Torus-equivariant CY$_3$ condition}: $T = (\CC^\times)^3$
preserves $\Omega$ iff $t_1 t_2 t_3 = 1$, equivalently
$\epsilon_1 + \epsilon_2 + \epsilon_3 = 0$.
\end{enumerate}
Under the Jordan triple presentation, (i), (ii), (iii) are three ways
of expressing the same CY$_3$ structure: the quiver-potential pair
$(Q, W)$ presents the non-commutative crepant resolution of $\CC^3$
at the symmetric position, and the torus action is then the isotropy.
The sentence ``The CY$_3$ condition is $t_1 t_2 t_3 = 1$'' should
read ``The $T$-equivariant CY$_3$ identity, consistent with the
algebraic and geometric CY$_3$ conditions, is $t_1 t_2 t_3 = 1$.''
\end{correction}

\subsection{Six-dimensional hCS: the $E_n$-class and Dolbeault
locality}
\label{subsec:6d-hCS}

\begin{flaw}[``holomorphic $E_3$-algebra'' is under-defined in the
treatise]
\label{flaw:hol-e3}
The treatise's subsubsection ``$E_3$-algebra structure of
Chern--Simons observables'' asserts:
\emph{``The observable algebra is neither $E_3$ (topological) nor
$E_1$ (just associative), but rather a holomorphic $E_3$-algebra,
i.e.\ a factorisation algebra on $\CC^3$ with holomorphic rather than
topological locality''.} The adjective ``holomorphic $E_3$'' is used
without a structural definition, and mixes two notions: (a) the
$\Ethree$-algebra governing the observables of a topological 3-manifold
and (b) the $\Ed$-factorisation-algebra structure given by Dolbeault
locality on a complex $d$-fold. Costello--Gwilliam 2017, Vol.~I
Thm.~5.3.3 gives (b) as holomorphically-local factorisation algebra.
\end{flaw}

\begin{correction}[Precise $\Ethree$-class of $\Obs_{\hCS}(\CC^3)$]
\label{cor:hol-e3}
Costello--Gwilliam Vol.~I Thm.~5.3.3 gives: a factorisation algebra on
$\RR^n$ with holomorphic locality (OPE singularities are poles in
Dolbeault coordinates) is equivalent to a vertex algebra on
$\CC^{n/2}$ when $n$ is even. For $\CC^3$ with $\hCS$, the precise
structure on observables is a factorisation algebra in
$\cD(\mathrm{Dolb}/\CC^3)$ which in \texttt{working\_notes.tex} and
\texttt{notes/platonic\_synthesis\_waves\_11\_through\_16.tex} is
recorded as ``$\Ethree$-$\HolFA(X)$ on the Calabi--Yau $d$-fold $X$''
(Thm.~\texttt{wn:thm:plat-hCS-quantum}), i.e.\ an $\Ethree$-algebra in
the symmetric-monoidal category of cochain complexes in Dolbeault
sheaves on $X$. Commutativity is $\pi_1(\mathrm{Conf}_2(\CC^3)) =
\pi_1(S^5) = 0$; associativity is Mayer--Vietoris on the
Fulton--MacPherson bar of $\CC^3$; deformation parameters are
$(\epsilon_1, \epsilon_2, \epsilon_3)$ with CY$_3$ identity and the
cohomological-degree shift $\hbar$ in the BV complex
(\texttt{working\_notes.tex} line 9--10, or the Platonic synthesis
Thm.~\texttt{wn:thm:plat-hCS-classical}).
\end{correction}

\begin{remark}[$\Ethree$ at $d=3$ and $\Phi$-output-scope]
The cache fact ``$d \geq 3$: $A$ is $\Eone$; $\Etwo$ lives on
$\cZ(\Rep(A))$, not on $A$'' applies on the curve side
(chiral image). On the CY side, the $\hCS$ observables carry an
$\Ed$-factorisation-algebra structure with $d = 3$, and its
boundary / specialisation to a curve produces the $\Eone$-chiral image
along that curve. This two-stage separation is stated exactly in
\texttt{notes/platonic\_synthesis\_waves\_11\_through\_16.tex}
Thm.~\texttt{wn:thm:plat-two-stage}: stage~1 gives $\Ed$-$\HolFA(X)$,
stage~2 specialises to $\Eone$-$\mathrm{ChirAlg}(C)$ over a
$(d{-}1)$-cycle $\Sigma_{d-1}$. The treatise's ``holomorphic
$E_3$-algebra'' is stage~1; its ``vertex algebra on a defect curve'' is
stage~2.
\end{remark}

\subsection{Conifold CoHA: Rap\v{c}\'ak--Soibelman--Yang--Zhao and the
toroidal double}
\label{subsec:conifold-toroidal}

\begin{flaw}[Conifold vertex-algebra image is scope-conflated]
\label{flaw:conifold-scope}
The treatise's conifold column asserts
\emph{``The positive-half CoHA of the conifold matches the positive
half of the quantum toroidal algebra
$U_{q_1, q_2}(\ggtwo)$''} (Rap\v{c}\'ak--Soibelman--Yang--Zhao 2020;
Kapranov--Vasserot 2018). The cohomological side
(Rap\v{c}\'ak--Soibelman--Yang--Zhao) and the $K$-theoretic side
(Kapranov--Vasserot) are different theorems on different cohomology
theories; the treatise's summary ``Drinfeld double Theorem (KV)''
identifies them. The two theorems are compatible but not the same;
the cohomological-to-$K$-theoretic comparison is via a formal-group
specialisation, not via identity.
\end{flaw}

\begin{correction}[Conifold: cohomological vs.\ $K$-theoretic scopes]
\label{cor:conifold-scope}
Rap\v{c}\'ak--Soibelman--Yang--Zhao 2020 (\texttt{arXiv:2001.10549})
Thm.~B: the cohomological CoHA of the conifold quiver at the
$T^2$-fixed locus is isomorphic to the positive half of a specific
affine shuffle algebra over $\FF[\epsilon_1, \epsilon_2]$ with kernel
$\omega_{\mathrm{con}}(z, w) = (z-w+\epsilon_1)(z-w+\epsilon_2)/
((z-w)(z-w+\epsilon_1+\epsilon_2))$. Kapranov--Vasserot 2018
(\texttt{arXiv:1802.07988}): the $K$-theoretic Hall algebra of the
conifold matches the positive half of quantum toroidal $\ggtwo$ at a
formal-group specialisation. Both are theorems; their mutual
compatibility follows from the expected shuffle-to-quantum-toroidal
map under $\epsilon \mapsto q - 1$ (\`a la Enriquez), and is not a
single-citation identity.
\end{correction}

\subsection{$K3 \times E$ CoHA: Davison hypothesis and its scope}
\label{subsec:davison-hypothesis}

\begin{flaw}[Davison CoHA stated in excess of hypothesis]
\label{flaw:davison-hypothesis}
The treatise states: \emph{``Davison 2017; Davison--Meinhardt 2020.
Scope hypothesis: $X$ a smooth proper Calabi--Yau three-fold admitting
a global critical-locus / potential presentation for its sheaf moduli
(not automatic for every compact CY$_3$; Davison's construction is
formulated for quiver-with-potential models and their direct geometric
counterparts, with the global critical-chart hypothesis load-bearing).
Under this hypothesis, there is a cohomological Hall algebra
$\CoHA^{\mathrm{BPS}}(X)$...''.} The scope is underspecified. The
Davison--Meinhardt 2020 (\emph{Invent.\ Math.} 221) integrality
framework applies to quivers with potential and certain global Joyce
d-critical structures; it does \emph{not} apply generally to smooth
proper CY$_3$ without a global critical-locus presentation. On
$K3 \times E$, no such presentation has been proved (only conjectured
via local-to-global gluing).
\end{flaw}

\begin{correction}[Davison scope on $K3 \times E$]
\label{cor:davison-k3e}
Davison 2017 (\texttt{arXiv:1601.02479}) and Davison--Meinhardt 2020
(\emph{Invent.\ Math.} 221) establish: (a) for any quiver with
potential $(Q, W)$, the critical cohomology $H^*(\cM(Q, d), \phi_W
\cdot \mathrm{IC})$ carries an associative CoHA structure; (b) a BPS
Lie subalgebra $\fg_{\mathrm{BPS}}$ of primitive elements exists, and
the full CoHA is its universal enveloping algebra; (c) the integration
map to motivic Hall algebra is compatible. These apply to smooth
proper CY$_3$ only under a global critical-locus /
d-critical-atlas hypothesis, which is \emph{not} known for
$K3 \times E$ in general. The $K3 \times E$ instantiation is therefore
recorded in \texttt{working\_notes.tex} line 9698 as conjectural:
``$\CoHA(K3 \times E)$ is $Y^+(\fg_{\Delta_5})$, the positive half
only (not the full BKM, by the Vol III cache discipline)''. The
treatise's row ``CoHA as associative: Theorem (Davison 2017)'' in the
$K3 \times E$ column should read ``conjecturally under global
critical-chart hypothesis''.
\end{correction}

\subsection{Oberdieck--Pixton scope: primitive K3 class only}
\label{subsec:op-scope}

\begin{flaw}[Oberdieck--Pixton scope is correctly noted but then
dropped]
\label{flaw:op-scope}
The treatise correctly scopes the Oberdieck--Pixton 2017 theorem at
the primitive K3 class (\texttt{arXiv:1706.10100}, \emph{Invent.\
Math.} 222 (2020)): ``\emph{Scope restriction: the theorem as stated
covers the primitive-class reduced DT partition function; the
imprimitive / multiple-cover contribution is not inside the stated
identity.}'' However, two paragraphs later it writes
\emph{``$\Omega(\gamma) = \mathrm{mult}_{\mathrm{BKM}}(\alpha_\gamma)$
for all $\gamma$''} and then
\emph{``$\dim \fg_{\mathrm{BPS}, \gamma}(K3 \times E) =
\mathrm{mult}_{\mathrm{BKM}}(\alpha_\gamma)$''}. The quantifier ``for
all $\gamma$'' exceeds the primitive-class scope of the Oberdieck--Pixton
theorem; non-primitive classes have not been reduced to the
Borcherds-product coefficient.
\end{flaw}

\begin{correction}[Primitive-class scope on
$\Omega(\gamma) = \mathrm{mult}_{\mathrm{BKM}}$]
\label{cor:primitive-only}
The identity $\Omega(\gamma) = \mathrm{mult}_{\mathrm{BKM}}(\alpha_\gamma)$
is proved only for primitive K3 classes (Oberdieck--Pixton 2017;
Pandharipande--Thomas 2014). For imprimitive classes, the
multiple-cover / stable-pairs corrections are not known to match the
Borcherds-product coefficients in closed form. The Vol III programme
records this as \texttt{working\_notes.tex} Thm.~5.2
``Oberdieck--Pixton 2017: primitive K3 class'' and Open Problem~7
``imprimitive class reduction to $\Phi_{10}^{-1}$''.
\end{correction}

\subsection{Mukai-group equivariant CoHA: speculative attribution}
\label{subsec:mukai-equivariant}

\begin{flaw}[Treatise presents $G \times E[N]$-equivariant CoHA as
programme-standard]
\label{flaw:mukai-equivariant}
The subsubsection ``Can we construct the positive-half CoHA...''
lists Nikulin 1980 and Mukai 1988 and then writes:
\emph{``For each Mukai group $G$, $G$-equivariant cohomology
$H^*_G(K3)$ is defined. Combined with $E$-translation subgroup
$E[N] \cong (\ZZ/N)^2$, one gets $G \times E[N]$-equivariant
cohomology of $K3 \times E$.''} The Mukai equivariant cohomology of
K3 is classical; the $G \times E[N]$-equivariant cohomology of
$K3 \times E$ is defined, but the subsequent ``$G \times E[N]$-equivariant
CoHA'' has not been constructed in the literature. Schiffmann--Vasserot
2013 constructs the Heisenberg-type cohomological Hall algebra of K3
via the Nakajima action; no equivariant extension to $K3 \times E$
appears in primary literature. The treatise flags ``\emph{Status of
construction: this has not been carried out in the literature to my
knowledge}'' but in summary-tabular form the row is left as
``Partial (Davison 2017)''.
\end{flaw}

\begin{correction}[Mukai-equivariant CoHA on $K3 \times E$ is open]
\label{cor:mukai-equivariant-open}
Scope-wise, the construction of a $G \times E[N]$-equivariant
cohomological Hall algebra of the stack of semistable sheaves on
$K3 \times E$ with $G$ a Mukai-finite-symplectic-automorphism group
and $E[N]$ the translation subgroup of the elliptic curve is \emph{open}
(no single-citation theorem or preprint exists as of 2026-04-22). The
programme's placeholder is recorded in Platonic Synthesis
(\texttt{notes/platonic\_synthesis\_waves\_11\_through\_16.tex}) under
``Residual frontier: Mordell--Weil-indexing conjecture'', where six
routes to $G(K3 \times E)$ are unified as $\mathrm{Aut}(K3) \times
\mathrm{SL}_2(\ZZ)$-orbits in $H_4(K3 \times E; \ZZ)$.
\end{correction}

\subsection{Six routes and the mis-framed ``six $\Phi$-applications''}
\label{subsec:six-routes}

\begin{flaw}[Treatise's closing paragraph blurs ``six routes'']
\label{flaw:six-routes}
The summary table closes with the row
\emph{``BPS Lie algebra: $\widehat{\fgl}_1$ / toroidal $\widehat{\widehat{\fgl}_2}$ /
$\fg_{\Delta_5}$?''}. The treatise does not explicitly claim ``six
$\Phi$-applications to $K3 \times E$'', but the enumeration of
``multiple candidate $\Sigma$ choices'' at the ``What would a chiral
BKM on $K3 \times E$ look like?'' subsection ($\Sigma = E$; $\Sigma =
\CC$; $\Sigma = $ genus-2 hyperelliptic in $K3$) invites the
mis-reading that the resulting chiral algebras are all the image of
$\Phi$ on the \emph{same} CY datum.
\end{flaw}

\begin{correction}[Six routes are six DIFFERENT constructions]
\label{cor:six-routes}
The cache-canonical formulation
(\texttt{appendices/first\_principles\_cache.md}): the six routes to
$G(K3 \times E)$ are six \emph{different} constructions --- Mukai
lattice, Igusa $\Phi_{10}$ via Gritsenko $\Delta_5$, BKM Borcherds
weight, K3 fibre-rank, Heisenberg--Mukai, Maulik--Okounkov
$R$-matrix --- \emph{not} six applications of $\Phi$ to one object.
The two-stage factorisation
(\texttt{notes/platonic\_synthesis\_waves\_11\_through\_16.tex}
Thm.~\texttt{wn:thm:plat-two-stage}) sharpens this: a single
$\cF_{K3 \times E} \in E_3\text{-}\HolFA(K3 \times E)$ is
$\Phi$-canonical at stage~1, and the zoo of chiral images arises from
stage~2 specialisations $\mathrm{Sp}_{\Sigma_2, C}$. Different
$(\Sigma_2, C)$ produce different $\Eone$-chiral shadows. The
treatise's choice-of-$\Sigma$ paragraph is directionally correct but
should cite the two-stage theorem as the frame.
\end{correction}

\subsection{Abelian 4D $\mathcal{N} = 2$ $U(1)$ with $\Omega$-background:
partition function}
\label{subsec:nekrasov-abelian}

\begin{flaw}[$Z^{\mathrm{inst}}$ on abelian 4D $\mathcal{N} = 2$ is
stated as pure $U(1)$, but formula misses $\epsilon$-dependence]
\label{flaw:nekrasov-u1}
The treatise writes:
\[
Z^{\mathrm{inst}}_\Omega(\epsilon_1, \epsilon_2; q)
= \sum_{n \geq 0} q^n \chi_T(\Hilb^n(\CC^2))
= \prod_{n \geq 1} \frac{1}{1 - q^n}.
\]
In Nekrasov 2003 (\texttt{arXiv:hep-th/0206161}), the abelian
(pure $U(1)$) instanton partition function is
$\eta(q)^{-1} = q^{-1/24} \prod_n (1 - q^n)^{-1}$ up to the vacuum
prefactor, but the $T$-equivariant $\chi_T(\Hilb^n(\CC^2))$ on
$\Omega$-background has explicit $\epsilon$-dependence via the
tangent weights at fixed Young-tableau loci; the reduction to
$\prod (1 - q^n)^{-1}$ is at the $\epsilon$-equivariant-Euler-number
level only (i.e.\ $\chi_T|_{\epsilon_1, \epsilon_2 \to 0}$). The
treatise's equality is therefore an identity between the
$\epsilon$-independent limit of the Hilbert-scheme Euler number and the
dedekind-eta character, not the fully equivariant Nekrasov partition
function.
\end{flaw}

\begin{correction}[Abelian Nekrasov function, made precise]
\label{cor:nekrasov-u1}
The fully equivariant $U(1)$ Nekrasov partition function with
$\Omega$-background weights $(\epsilon_1, \epsilon_2)$ is
\[
Z^{\mathrm{inst}}_{\Omega, U(1)}(\epsilon_1, \epsilon_2; q)
= \sum_{\mu} q^{|\mu|}
\prod_{s \in \mu} \frac{1}{a_\mu(s) \cdot \epsilon_1 - (l_\mu(s) + 1)
\cdot \epsilon_2) \cdot ((a_\mu(s) + 1) \cdot \epsilon_1 -
l_\mu(s) \cdot \epsilon_2)},
\]
with the sum over Young tableaux $\mu$ and $a_\mu, l_\mu$ the arm and
leg lengths. The $\epsilon_1, \epsilon_2 \to 0$ limit reduces to
$q^{-1/24} / \eta(q)$ (Nakajima--Yoshioka 2003, Thm.~2.1). The
treatise's $\prod (1 - q^n)^{-1}$ formula is the
$\epsilon$-independent-Euler-character limit, which is correct as a
combinatorial identity but mis-labelled as the equivariant
partition function.
\end{correction}

\subsection{Bare $\kappa$ and other antipattern violations}
\label{subsec:antipatterns}

The treatise uses no bare $\kappa$ symbol (it does not perform
$\kappa$-arithmetic at all), so AP113 is not directly violated.
However, several Pattern-adjacent infelicities:
\begin{itemize}
\item The row ``Identified with VOA'' in the $\CC^3$ column violates
Pattern 273 ($\Phi$ functor vs.\ object-level correspondence):
$\Phi(\CC^3) = \Winf[\lambda]$ is a statement at the image-of-$\Phi$
level only after specialisation, not an object-level identification
with $\CoHA(\CC^3) = Y^+(\gghone)$.
\item The phrase ``identified with $\Winf$-affine Yangian'' in the
Tsymbaliuk 2017 theorem statement conflates the associative affine
Yangian $Y(\gghone)$ with the vertex algebra $\Winf[\lambda]$; these
are different algebraic classes (associative vs.\ chiral), and the
Tsymbaliuk theorem concerns the former, not the latter.
\item The treatise tags
``$\CoHA(\CC^3) \otimes_\FF \FF(\!(z)\!) \cong \mathrm{Sh}$'' as
$\mathtt{ClaimStatusTheorem}$, which is correct (Schiffmann--Vasserot
2013 Thm.~1.1), but the subsequent ``Drinfeld double is isomorphic,
as a topological Hopf algebra, to the affine Yangian of $\gghone$ in
the Drinfeld-currents presentation with structure function
$\varphi(u) = (u + \epsilon_1)(u + \epsilon_2)(u + \epsilon_3)/u^3$''
is a Tsymbaliuk 2017 theorem that should be cited with the full
convention-fixing statement: the positive-half generators are
identified on a specific branch of a square-root of the structure
function (Tsymbaliuk 2017 Thm.~1.1 is more careful on this branch
choice than the treatise paraphrases).
\item The treatise's claim
``\emph{A factorisation algebra on $\CC$ with holomorphic locality is
equivalent to a vertex algebra (Costello--Gwilliam 2017 Vol I
Thm.~5.3.3)}'' is essentially correct but the equivalence is
one-directional up to grading shift unless the factorisation algebra
is locally constant in a precise sense; Costello--Gwilliam's theorem
is technically precise on this.
\end{itemize}

\section{Consolidated correct dictionary}
\label{sec:dictionary}

Below is the dictionary in the form the treatise should carry once
the corrections of Section~\ref{sec:inconsistency} are applied. Each
entry states the CoHA-side object, the $\Winf$-side avatar, and the
nature of the map between them.

\begin{center}
\renewcommand{\arraystretch}{1.35}
\begin{tabular}{p{5.2cm}p{4.8cm}p{4.6cm}}
\toprule
\textbf{CoHA side ($Y^+(\gghone)$)} &
\textbf{Affine-Yangian side ($Y(\gghone)$)} &
\textbf{$\Winf$ side ($\Winf[\lambda]$)} \\
\midrule
$\CoHA(\CC^3) = Y^+(\gghone)$: associative
($\Eone$), shuffle presentation with kernel
$\omega(z, w) = \prod_i (z - w - \epsilon_i)/(z-w)^3$,
generated by power sums $p_k$ (Schiffmann--Vasserot 2013 Thm.~1.1;
Tsymbaliuk 2017 Thm.~1.1). &
Drinfeld double $Y = Y^+ \bowtie Y^0 \bowtie Y^-$: topological
Hopf algebra in Drinfeld-currents presentation with structure
function $\varphi(u) = \prod_i (u + \epsilon_i)/u^3$
(Tsymbaliuk 2017 Thm.~1.1). &
$\Winf[\lambda]$: vertex algebra on a curve, central charge $c$,
generating fields at each conformal dimension $h = 1, 2, 3, \ldots$
(Kac--Radul 1993; Linshaw 2011). \\
Associativity class: $\Eone$ (no OPE; associative composition of
correspondences). &
Associativity class: $\Eone$ (topological Hopf, not chiral). &
Associativity class: $\Etwo$ (chiral algebra on a curve; OPE
with singularities in $z - w$). \\
No conformal structure (CoHA carries only a grading by dimension
vector). &
Evaluation $\ev_z: Y \to \mathrm{End}$ on the spectral line (Costello--Yagi
2018; Maulik--Okounkov 2013). &
Full vertex-algebra conformal structure; Virasoro--$\Winf$ generators;
$\cN$-ality. \\
\bottomrule
\end{tabular}
\end{center}

\noindent
The three columns relate by:
\[
\CoHA(\CC^3) \;=\; Y^+(\gghone)
\xhookrightarrow{\text{Drinfeld double}}
Y(\gghone)
\xrightarrow[\ev_\lambda]{\text{state-field on vacuum}}
\mathrm{End}(\Winf[\lambda])\text{-vacuum}.
\]
The last arrow is a representation, not an isomorphism of
algebras: $Y(\gghone)$ acts on the vacuum module of $\Winf[\lambda]$
via the Gaiotto--Rap\v{c}\'ak 2017 construction, and the image is the
endomorphism vertex algebra of that module, which is precisely
$\Winf[\lambda]$. In particular,
\[
\CoHA(\CC^3) \;\neq\; \Winf,
\quad \text{and}\quad
Y(\gghone) \;\neq\; \Winf[\lambda] \text{ as an object (one is
associative Hopf, the other is chiral).}
\]

\subsection{Six-dimensional hCS on $\CC^3$: the $\Ethree$-factorisation
vantage}
\label{subsec:6d-hcs-vantage}

The relation $\CoHA(\CC^3) = Y^+(\gghone)$ versus the conjectural
$\Phi(\CC^3) \sim \Winf[\lambda]$ on a specified curve has a precise
origin in the six-dimensional holomorphic Chern--Simons theory on
$\CC^3$ with $\Omega$-background. The construction takes place at two
vertically-coupled stages:
\begin{enumerate}[label=\textbf{Stage \arabic*.}]
\item \emph{Canonical holomorphic $\Ethree$-factorisation algebra on
$\CC^3$.} The hCS observables $\Obs_{\hCS}(\CC^3)$ carry an
$\Ethree$-algebra structure in $\cD(\mathrm{Dolb}/\CC^3)$
(Costello--Gwilliam 2017 Vol.~I Thm.~5.3.3; Platonic Synthesis
Thm.~\texttt{wn:thm:plat-hCS-quantum}). Commutativity is
$\pi_1(\mathrm{Conf}_2(\CC^3)) = \pi_1(S^5) = 0$; associativity is
Mayer--Vietoris on Fulton--MacPherson configuration space.
\item \emph{Specialisation to a defect line $\Sigma = \CC_z \subset
\CC^3$.} Factorisation homology of $\cF$ over the transverse
$(d-1) = 2$-cycle $\Sigma_2 \subset \CC^3$, restricted to the defect
line $C \cong \CC_z$, produces an $\Eone$-chiral algebra on $C$. This
is stage~2 of Thm.~\texttt{wn:thm:plat-two-stage}.
\end{enumerate}
Under the $\Omega$-background with weights $(\epsilon_1, \epsilon_2,
\epsilon_3)$, $\epsilon_1 + \epsilon_2 + \epsilon_3 = 0$, the stage-2
output is $\Winf[\lambda]$ with $\lambda$ determined by the
transverse weights (Costello 2013 \texttt{arXiv:1303.2632}; Gaiotto--Rap\v{c}\'ak
2017; Costello--Paquette 2020 \texttt{arXiv:2009.04834}). The
CoHA-to-$\Winf$ map is then the composite
\[
\CoHA(\CC^3)
\xrightarrow{\mathrm{SV\ shuffle}}
Y^+(\gghone)
\hookrightarrow Y(\gghone)
\xrightarrow{\mathrm{Sp}_{\Sigma_2, C}}
\Winf[\lambda],
\]
where $\mathrm{Sp}_{\Sigma_2, C}$ is the stage-2 specialisation.
The treatise's treatment is correct on stage~1 (though phrased as
``holomorphic $E_3$-algebra'' without the $\cD(\mathrm{Dolb})$ target)
and ambiguous on stage~2 (where it conflates the representation
$Y \to \mathrm{End}(\Winf)$ with an isomorphism $Y \cong \Winf$).

\subsection{The 6d hCS on a larger space: a structural conjecture
about the treatise's hypothetical ``$W_{1+\infty}$ direct image''}
\label{subsec:6d-hcs-larger-space}

Is there a 6d-hCS ambient space on which the boundary of the
holomorphic Chern--Simons theory is \emph{directly} $\Winf$?
Directionally, the programme records two genuine routes:
\begin{enumerate}[label=\textup{(\roman*)}]
\item \emph{Class-$\mathcal{S}$ route.} Compactify the 6d $(2,0)$-theory
on a Riemann surface with sufficient punctures; for the
24-punctured sphere $\Sigma_{0, 24}$, Chacaltana--Distler 2010
(Table~3 row~1) gives $c_{4d} = (5n - 13)/6 = 107/6$ at $n = 24$
(Platonic Synthesis retraction~6). The 4d $\mathcal{N} = 2$
$\Omega$-background construction then produces a 2d chiral algebra
which can be identified, at the tame-limit point, with a
Virasoro/$W$-extension. This is \emph{not} $\Winf$; it is a
Virasoro-extended vertex algebra with central charge $107/6$.
\item \emph{Gaiotto--Rap\v{c}\'ak Y-algebra route.} The
Gaiotto--Rap\v{c}\'ak $Y$-algebras $Y_{L, M, N}[\psi]$ are
parametrised vertex algebras on the junction of three-Chan--Paton
branes; $Y_{0, 0, N}[\psi]$ is the $\Winf$-sector at truncation
$N$. The ambient space is \emph{not} a higher-dimensional hCS
target; it is a 5-brane web reduction.
\end{enumerate}
The structural conjecture that ``there is a 6d hCS on a \emph{larger}
space whose boundary is $\Winf$'' is directionally consistent with
what Costello--Paquette 2020 call the holographic dual of
4d~$\mathcal{N}=2$~$\Omega$-background, but neither Costello--Paquette
nor Costello 2013 produces such a higher-dimensional hCS target as
primary source. Treating this as an \emph{open structural question}
rather than a claim: \emph{Is there a canonical 6-manifold $X^{(6)}
\supset \CC^3$ on which holomorphic Chern--Simons has, as its
$\Eone$-chiral shadow along a distinguished curve, the algebra
$\Winf[\lambda]$ directly, without the specialisation
$\mathrm{Sp}_{\Sigma_2, C}$ of Thm.~\texttt{wn:thm:plat-two-stage}?}
The Vol III programme records this as the residual-frontier item
``elliptic-surface specialisation $(\Sigma_2, C) = (\mathcal{E},
\PP^1)$: structural frontier; if Mordell--Weil-indexing conjecture
holds, the six routes unify''.

\section{Cross-consistency verdict}
\label{sec:cross-consistency}

\begin{enumerate}
\item \textbf{Consistent with cache:} Jordan-triple Jacobi computation,
Schiffmann--Vasserot Thm.~1.1, Tsymbaliuk 2017 Drinfeld-double
isomorphism, Szendroi 2008 conifold NCCR, Negut 2015 conifold shuffle
kernel, Davison--Meinhardt 2020 associativity, Oberdieck--Pixton 2017
primitive-class identification, Gritsenko--Nikulin 1998 Borcherds
product for $\Delta_5$, Costello--Gwilliam 2017 factorisation-algebra
framework.

\item \textbf{Inconsistent with cache:} (i) Implicit identification
$\CoHA(\CC^3) = \Winf$ via ``identified with VOA'' row
(Section~\ref{subsec:coha-vs-winf}); (ii) CY$_3$ torus-equivariant
identity conflated with algebraic CY$_3$ condition
(Section~\ref{subsec:cy3-identity}); (iii) Oberdieck--Pixton scope
dropped at ``$\Omega(\gamma) = \mathrm{mult}_{\mathrm{BKM}}$ for all
$\gamma$'' (Section~\ref{subsec:op-scope}); (iv) Davison CoHA stated
in excess of scope for $K3 \times E$
(Section~\ref{subsec:davison-hypothesis}); (v) Mukai-group equivariant
CoHA presented as programme-standard without primary literature
(Section~\ref{subsec:mukai-equivariant}); (vi) Abelian Nekrasov
partition function formula confuses equivariant and Euler-character
limits (Section~\ref{subsec:nekrasov-abelian}); (vii)
``Holomorphic $E_3$-algebra'' used without structural definition
(Section~\ref{subsec:6d-hCS}); (viii) Conifold RSYZ / KV attributions
conflated across cohomology theories
(Section~\ref{subsec:conifold-toroidal}).

\item \textbf{Inconsistencies vs.\ \texttt{working\_notes.tex}:}
The treatise's $\CC^3$ summary column stands at the surface of
\texttt{working\_notes.tex} \S\ref{sec:c3-rosetta} (``$\C^3$ Rosetta
Stone''), which already carries the programme-standard
formulation: $\CoHA(\CC^3) = Y^+(\gghone) \hookrightarrow Y(\gghone)
\to \Rep^{\Etwo}(\Winf)$ via the Drinfeld center
(Thm.~\texttt{wn:thm:c3-drinfeld-center}). The five routes to $\Winf$
(\texttt{working\_notes.tex} line 1665) are:
(a) critical CoHA;
(b) crystal melting;
(c) shuffle algebra;
(d) factorisation envelope;
(e) 5d Chern--Simons.
All five produce $\Winf$ at $N = 1$; none equates $\CoHA$ with
$\Winf$.

\item \textbf{Consistent with Platonic Synthesis Waves 11--16:} The
two-stage factorisation (Thm.~\texttt{wn:thm:plat-two-stage}) and
the CoHA/Miki triality (Thm.~\texttt{wn:thm:plat-Miki-S3}, which
explicitly writes ``The cache rule $\CoHA(\CC^3) = Y^+ \neq \Winf$
stands: $Y^+$ is associative cohomological-Hecke; $\Winf$ is its
evaluation slice'') together provide the correct frame. The
treatise would benefit from adopting this language.

\item \textbf{Directional verdict.} Directionally, the treatise is
correct end-to-end: it identifies the right CoHA objects, the
correct Drinfeld-double presentation, the correct Borcherds-product
match at $K3 \times E$, and the correct scope of the $K3 \times E$
open problems. The inconsistencies are structural/labelling
inconsistencies, not first-principles errors.
\end{enumerate}

\section{Open questions sharpened by the audit}
\label{sec:open-questions}

\begin{enumerate}
\item \textbf{Global critical-locus hypothesis on $K3 \times E$.}
Does $K3 \times E$ carry a global d-critical /
critical-chart-atlas structure in the sense of Joyce required for
the Davison--Meinhardt 2020 theory to apply on the nose? If so, at
what level of resolution (\v{C}ech, \'etale, orbifold)?
\item \textbf{Mukai equivariant CoHA.} Construct, as a theorem,
the $G \times E[N]$-equivariant cohomological Hall algebra
$\CoHA^{G_N}(K3 \times E)$ for programme-core $N \in \{1, 2, 3, 4, 6\}$,
and compare its positive half to the Drinfeld-double positive half
of a Yangian associated to $\fg_{\Delta_5}$.
\item \textbf{Imprimitive-class BKM match.} Extend Oberdieck--Pixton
2017 beyond primitive K3 class; match imprimitive-class DT
invariants to Gritsenko--Nikulin Borcherds-product coefficients at
$\Delta_5$.
\item \textbf{6d hCS containment $\supset \CC^3$ with $\Winf$ direct
shadow.} Identify a canonical 6-manifold $X^{(6)} \supset \CC^3$
(or $X^{(6)}$ genuinely complex 3-fold, not $\CC^3$) on which the
holomorphic Chern--Simons boundary is $\Winf[\lambda]$ directly,
without specialisation.
\end{enumerate}

\appendix

\section{Rectification log}
\label{app:rectification-log}

Five attack-heal cycles; line references are to
\texttt{notes/CoHA\_to\_W\_infty\_treatise.tex} as read on 2026-04-22.

\begin{enumerate}
\item \textbf{Cycle 1 (lines 155--220; Tsymbaliuk + Kac--Radul block).}
ATTACK: ``Drinfeld double of $Y^+$ is $Y(\gghone)$ which is affine
Yangian'' is conflated with ``$Y$ is $\Winf$'' in the summary row.
HEAL: Correction~\ref{cor:coha-winf} and dictionary
Section~\ref{sec:dictionary}: $Y(\gghone)$ acts on $\Winf[\lambda]$
via $\ev_\lambda$; the map is a representation, not an isomorphism.

\item \textbf{Cycle 2 (lines 195--220; $\lambda$-parametrisation).}
ATTACK: Two $\lambda$-parameters coexist; the ``$\lambda = -1$ slip''
is framed as a single-point conflation. HEAL: Correction
\ref{cor:lambda-slices}: two distinguished slices exist after the
CY$_3$ identity, the self-dual $\epsilon_3 = 0$ (Heisenberg) and the
Miki-symmetric unrefined $\epsilon_1 = \epsilon_2$.

\item \textbf{Cycle 3 (lines 440--625; $K3 \times E$ block).}
ATTACK: Davison scope and Mukai-equivariant CoHA construction are
claimed in excess of primary literature. HEAL: Corrections
\ref{cor:davison-k3e} and \ref{cor:mukai-equivariant-open}: Davison
applies under global critical-chart hypothesis only; Mukai-equivariant
extension is open.

\item \textbf{Cycle 4 (lines 481--531; OP dimension identity).}
ATTACK: ``for all $\gamma$'' exceeds Oberdieck--Pixton primitive-class
scope. HEAL: Correction~\ref{cor:primitive-only}; the cache-canonical
identification is on primitive class only.

\item \textbf{Cycle 5 (lines 260--330; hCS block).}
ATTACK: ``Holomorphic $E_3$-algebra'' is used without structural
definition; summary row ``E_3-CS Costello--Francis--Gwilliam
deformation: Applicable'' refers to a non-primary source. HEAL:
Correction~\ref{cor:hol-e3}: the precise structure is
$\Ethree$-algebra in $\cD(\mathrm{Dolb}/\CC^3)$ per Costello--Gwilliam
2017 and Platonic Synthesis
Thm.~\texttt{wn:thm:plat-hCS-quantum}; the Costello--Francis--Gwilliam
reference is a pending-literature item, already flagged in the
treatise's ``Scope note''.
\end{enumerate}

\end{document}
